\documentclass[aspectratio=169,11pt]{beamer}

\usepackage{iftex}
\ifXeTeX
  \usepackage{fontspec}
  \usepackage{xeCJK}
  \setmainfont{Times New Roman}
  \setCJKmainfont{SimSun}
\fi
\usepackage{booktabs}
\usepackage{array}
\usepackage{siunitx}
\usepackage{xurl}

\usetheme{Madrid}
\usecolortheme{default}
\setbeamertemplate{navigation symbols}{}
\setbeamertemplate{footline}[frame number]

\newcommand{\movetag}[2]{\textbf{Move #1:} #2}
\newcommand{\steptag}[2]{\textbf{Step #1:} #2}
\newcommand{\claim}[1]{\textcolor{blue!65!black}{#1}}
\newcommand{\gaptext}[1]{\textcolor{red!65!black}{#1}}
\newcommand{\solutiontext}[1]{\textcolor{green!45!black}{#1}}
\newcommand{\src}[1]{\textsuperscript{\scriptsize[\texttt{\detokenize{#1}}]}}
\newcolumntype{L}[1]{>{\raggedright\arraybackslash}p{#1}}

\title[MICE Introduction Map]{Introduction 的 MICE 结构映射}
\subtitle{从当前情境到本文 surrogate workflow}
\author{Linan Jiang}
\institute{Aalto University}
\date{7 April 2026}

\begin{document}

\begin{frame}
  \titlepage
\end{frame}

\begin{frame}{MICE 模型：工程硕士论文引言的四个 Move}
\small
\begin{tabular}{L{0.18\linewidth}L{0.34\linewidth}L{0.38\linewidth}}
\toprule
Move & 功能 & 本文对应问题 \\
\midrule
Move 1 & Current situation & 为什么船用多燃料发动机中的 pilot spray 值得研究？ \\
Move 2 & Problems & 现有试验、CFD、经验模型和图像测量缺少什么？ \\
Move 3 & New approach & 为什么从 Mie 视频学习 surrogate 是合理路线？ \\
Move 4 & Your solution & 本论文具体做什么、怎么做、做到什么范围？ \\
\bottomrule
\end{tabular}

\vspace{0.7em}
\scriptsize
MICE 强调 general-to-specific 和 problem-solution 顺序；本 slide deck 用它检查 introduction 的论证链。
\end{frame}

\begin{frame}{Introduction 的总路线}
\footnotesize
\begin{enumerate}
  \item \claim{行业背景}：航运减排和替代燃料推动多燃料船用发动机发展。
  \item \gaptext{第一层问题}：替代燃料进入发动机后，可靠点火并不会自动解决。
  \item \claim{技术背景}：pilot diesel injection 仍是点燃低反应性主燃料的关键接口。
  \item \gaptext{第二层问题}：pilot spray 的穿深控制决定点火核位置、主燃料混合与燃烧效率。
  \item \gaptext{边界风险}：若这一空间控制失效，同一个 pilot spray 也会带来撞壁、润滑油稀释和混合恶化风险。
  \item \claim{已有基础}：试验、CFD、经验穿深模型、Mie 成像和数据驱动 surrogate 都提供了部分入口。
  \item \gaptext{方法缺口}：这些入口难以同时处理短脉冲、右删失、多孔配准和不确定性筛查。
  \item \solutiontext{新路线}：用 Mie 高速成像数据构造 reduced-order penetration targets，并训练不确定性感知 surrogate。
\end{enumerate}
\end{frame}

\begin{frame}{Move 1/2 的交叉 loop}
\small
\begin{tabular}{L{0.27\linewidth}L{0.30\linewidth}L{0.34\linewidth}}
\toprule
回合 & Move 1：建立背景/已有基础 & Move 2：立刻暴露限制 \\
\midrule
能源转型 & 替代燃料船用发动机变得重要。 & 低反应性燃料仍需要可靠点火。 \\
pilot spray & pilot diesel 是关键点火接口。 & 穿深控制决定点火核位置和主燃料混合；失控后才转化为撞壁边界风险。 \\
工程评估 & 试验、CFD、经验模型和 Mie 成像已有基础。 & 成本、速度、瞬态适用性和测量定义各有弱点。 \\
数据路线 & 神经 surrogate 已能预测喷雾宏观特性。 & 缺少短脉冲+右删失+多孔 plume+不确定性输出的组合。 \\
\bottomrule
\end{tabular}

\vspace{0.7em}
\footnotesize
这样写比单次 ``background $\rightarrow$ gap'' 更有张力：每引入一层已有知识，就说明为什么它还不足以完成本文任务。
\end{frame}

\section{Move 1}

\begin{frame}{Move 1-1：Making a centrality claim}
\small
\movetag{1}{Current situation}

\vspace{0.5em}
\steptag{1}{Making a centrality claim}

\vspace{0.8em}
\begin{block}{本文论证}
国际航运正同时面临温室气体减排、燃料多样化和运行可靠性压力；IMO 2023 战略和行业订单趋势使替代燃料推进系统成为现实工程问题 \src{imo2023ghgstrategy}。低压天然气、甲醇和氨等路线使双燃料/多燃料船用发动机成为重要开发方向 \src{classnk2025alternativefuels}。
\end{block}

\begin{itemize}
  \item 作用：把研究放在船用能源转型的大背景中。
  \item 写法：先说明 application 为什么重要，再引出它内部尚未解决的小问题。
  \item 引言位置：引言第 1 段。
\end{itemize}
\end{frame}

\begin{frame}{Move 1-2：Making topic generalizations}
\small
\movetag{1}{Current situation}

\vspace{0.5em}
\steptag{2}{Making topic generalizations}

\vspace{0.8em}
\begin{block}{本文论证}
双燃料和多燃料发动机不仅需要替代燃料供应，也需要可靠点火。pilot diesel injection 是点燃低反应性主燃料的关键接口；其穿深、锥角和液相包络决定点火核位置、主燃料混合和火焰传播，因此直接影响发动机方案能否稳定工作 \src{wingd_xdf_manual_2021; classnk_methanol_2024}。
\end{block}

\begin{itemize}
  \item 作用：把大背景收窄到 pilot spray。
  \item loop 中的下一句：若 pilot spray 的空间控制失效，它才进一步成为撞壁风险来源。
  \item 引言位置：引言第 2--4 段。
\end{itemize}
\end{frame}

\begin{frame}{Move 1-3：Reviewing previous research and solutions}
\small
\movetag{1}{Current situation}

\vspace{0.5em}
\steptag{3}{Reviewing previous research and solutions}

\vspace{0.8em}
\begin{block}{本文论证}
pilot spray 空间发展及其撞壁边界通常通过发动机试验、CFD/喷雾燃烧仿真、经验穿深模型和光学喷雾实验评估；高速 Mie 成像能提供受控条件下的液相包络证据，数据驱动 surrogate 也已显示出预测喷雾宏观特性的潜力 \src{Linne2013SprayImaging; HiroyasuArai1990; NaberSiebers1996; VazKarathanasis2026; liu2020sprayreview}。
\end{block}

\begin{itemize}
  \item 作用：承认领域已有成熟基础，避免把问题写成“没人研究过喷雾”。
  \item loop 中的下一句：然而，这些入口单独使用都不足以形成快速、可重复、带不确定性的筛查模型。
  \item 引言位置：引言第 5--9 段。
\end{itemize}
\end{frame}

\section{Move 2}

\begin{frame}{Move 2-1：Indicating weaknesses}
\small
\movetag{2}{Problems}

\vspace{0.5em}
\steptag{1}{Indicating weaknesses}

\vspace{0.8em}
\begin{itemize}
  \item 替代燃料并不自动解决点火问题，因此 pilot diesel injection 仍不可忽略。
  \item pilot spray 的空间发展首先控制点火核位置、主燃料混合和燃烧效率。
  \item 若这一空间控制失效，液相发展过远才会带来壁面润湿、润滑油稀释和混合恶化。
  \item 发动机试验接近真实工况，但昂贵、慢，难以隔离单一设计因素。
  \item 高保真 CFD 信息丰富，但一次 reacting diesel-spray LES 可消耗约 48,000 CPU core-hours \src{CurtisNiemeyerSung2017}。
  \item 经验模型对短脉冲 SOI/EOI 瞬态不够稳健；光学边界又具有方法依赖性 \src{ZhouLiYi2021; ZhouLiLaiWei2019; ecn_liquid_penetration; ecn_liquid_penetration_accessory}。
\end{itemize}
\end{frame}

\begin{frame}{Move 2-2：Identifying a gap}
\small
\movetag{2}{Problems}

\vspace{0.5em}
\steptag{2}{Identifying a gap}

\vspace{0.8em}
\begin{block}{本文 gap statement}
缺少一种位于光学实验归档和高成本验证之间的 intermediate-layer predictor：它应能快速、可重复地预测短脉冲预喷柴油液相包络随时间的演化，并同时处理右删失、多孔 plume 配准和不确定性输出。
\end{block}

\begin{itemize}
  \item 不是缺少喷雾研究本身。
  \item 缺的是把 Mie 视频证据转成可查询、可审计、可用于预喷控制比较和撞壁排序的模型层。
  \item 目标不是替代 CFD 或发动机试验，而是在它们之前收窄候选空间。
\end{itemize}
\end{frame}

\section{Move 3}

\begin{frame}{Move 3-1：Introducing a potential solution}
\small
\movetag{3}{Proposing a new approach}

\vspace{0.5em}
\steptag{1}{Introducing a potential solution}

\vspace{0.8em}
\begin{block}{本文新路线}
直接从大规模 Mie-scattering optical spray data 中学习一个 condition-parametric surrogate，把测得的 penetration trajectories 转换成可快速查询的预测模型。
\end{block}

\begin{itemize}
  \item Mie 视频已经捕获液相包络行为。
  \item Reduced-order fitting 把原始观测转换为稳定学习目标。
  \item Surrogate 模型位于 raw measurement 和 physics-based simulation 之间。
\end{itemize}
\end{frame}

\begin{frame}{Move 3-2：Justifying the approach}
\small
\movetag{3}{Proposing a new approach}

\vspace{0.5em}
\steptag{2}{Justifying the approach}

\vspace{0.8em}
\begin{itemize}
  \item \solutiontext{保留实验保真度}：训练目标来自实际 Mie 成像和穿深轨迹。
  \item \solutiontext{提高筛选速度}：避免每个候选工况都进入 CFD 或发动机试验。
  \item \solutiontext{适应数据扩展}：条件到轨迹参数的映射可随数据集更新。
  \item \solutiontext{可解释和可审计}：标定、边界定义、轨迹拟合和不确定性回归都有中间产物。
  \item \solutiontext{符合科学图像处理实践}：Fiji/ImageJ、scikit-image、OpenPIV、PIVlab 等生态强调脚本化、可检查、可复用流程 \src{Schindelin2012Fiji; VanDerWalt2014ScikitImage; OpenPIV; ThielickeSonntag2021PIVlab}。
\end{itemize}
\end{frame}

\section{Move 4}

\begin{frame}{Move 4-1：Aim and tasks}
\footnotesize
\movetag{4}{Your solution}

\vspace{0.5em}
\steptag{1}{Announcing aim and tasks}

\vspace{0.5em}
\begin{block}{本文目标}
把大规模 Mie 散射喷雾视频归档转化为位于光学实验归档与高成本验证之间的早期工程筛查层。
\end{block}

\vspace{0.3em}
\begin{enumerate}
  \item 提取逐 plume 几何观测量。
  \item 构造删失感知 penetration targets。
  \item 训练不确定性感知 surrogate。
  \item 连接预喷控制评估与概率性 wall-impingement screening proxy。
\end{enumerate}
\end{frame}

\begin{frame}{Move 4-2：Methods / testbed}
\footnotesize
\movetag{4}{Your solution}

\vspace{0.5em}
\steptag{2}{Methods}

\vspace{0.5em}
\begin{enumerate}
  \item OSCC 室温、非反应、顶视高速 Mie 记录。
  \item 人工喷嘴标定：中心、半径、角向参考和像素尺度。
  \item CUDA 图像链：鲁棒预处理、Sobel、极坐标/FFT 配准、plume 条带提取。
  \item 轨迹链：CDF 穿深、边界度量、清洗、降阶拟合和右删失处理。
  \item 模型链：Stage~1 MSE、Stage~2 heteroscedastic NLL、Stage~3 raw-CDF/KD refinement。
\end{enumerate}
\end{frame}

\begin{frame}{Move 4-3：Thesis scope}
\small
\movetag{4}{Your solution}

\vspace{0.5em}
\steptag{3}{Thesis scope}

\vspace{0.8em}
\begin{itemize}
  \item 范围内：受控加压舱、室温、非反应、非蒸发 pilot diesel spray。
  \item 范围内：光学布置下观察到的宏观液相包络几何及其时间演化。
  \item 范围外：燃烧化学、蒸发、双燃料点火耦合、真实热壁面相互作用、真实发动机 wall-film truth。
  \item 模型边界：主要面向数据覆盖范围内的插值预测；明显 OOD 外推不作同等级保证。
\end{itemize}
\end{frame}

\begin{frame}{Move 4-4：Main outcome}
\footnotesize
\movetag{4}{Your solution}

\vspace{0.5em}
\steptag{4}{Describing the main outcome}

\vspace{0.5em}
\begin{block}{精简结果}
在外部清洗评估协议下，三阶段过渡区缩放 anchor-off Stage~3 surrogate 在 795{,}561 个预测点和 33{,}845 条轨迹上达到四次重复训练平均 RMSE~\SI{6.094}{mm}，并给出 0.783/0.982 的保守 1$\sigma$/2$\sigma$ 覆盖率。
\end{block}

\vspace{0.4em}
\begin{itemize}
  \item 较早的原始留出切分结果 RMSE~\SI{6.68}{mm}（21{,}470 条轨迹）保留为 legacy split 结果，而非 headline claim。
  \item 输出服务于预喷控制比较和可排序的概率性筛查，而不是二元通过/不通过结论。
  \item 结果讨论留给 Results 和 Discussion；Introduction 只给定位和可信度信号。
\end{itemize}
\end{frame}

\begin{frame}{Move 4-5：Thesis structure}
\scriptsize
\movetag{4}{Your solution}

\vspace{0.5em}
\steptag{5}{Overview of thesis structure}

\vspace{0.6em}
\begin{tabular}{L{0.29\linewidth}L{0.58\linewidth}}
\toprule
章节 & 作用 \\
\midrule
文献综述与技术依据 & 放置于柴油喷雾成像、光学诊断和 surrogate 建模背景下。 \\
实验数据与图像处理工作流 & OSCC 数据、人工标定、CUDA 预处理、配准、条带提取和 CDF/边界度量。 \\
轨迹目标、Surrogate 建模与概率性筛查 & 轨迹清洗、降阶拟合、特征设计、Stage~1--3 和碰壁筛查 proxy。 \\
结果 & 报告总体精度、campaign 差异、消融和原型筛查输出。 \\
讨论 & 解释 Stage~3 可信度、已验证内容和适用边界。 \\
结论 & 总结已验证结果、原型边界和下一步开发重点。 \\
\bottomrule
\end{tabular}
\end{frame}

\begin{frame}{Move/Step 检查表}
\scriptsize
\begin{tabular}{L{0.19\linewidth}L{0.30\linewidth}L{0.40\linewidth}}
\toprule
MICE 位置 & Introduction 小节 & 核心论证 \\
\midrule
Move 1-1 & 引言第 1 段 & 航运减排和替代燃料使多燃料发动机问题重要。 \\
Move 2-1 & 引言第 2 段 & 替代燃料仍需要可靠点火，问题收窄到 pilot injection。 \\
Move 1-2 & 引言第 3 段 & pilot diesel 是点燃低反应性主燃料的关键接口，且其空间发展控制点火核位置和主燃料燃烧。 \\
Move 2-1 & 引言第 4 段 & 当 pilot spray 空间控制失效时，才进一步表现为撞壁、润滑油稀释和混合恶化风险。 \\
Move 1-3 & 引言第 5/7/9 段 & 试验、CFD、经验模型、Mie 成像和数据 surrogate 已有基础。 \\
Move 2-1/2-2 & 引言第 6/8/10 段 & 仍缺少处理短脉冲、右删失、多孔配准和不确定性的中间层预测器。 \\
Move 3-1 & Problem statement 后段 & 从 Mie 数据学习 surrogate。 \\
Move 3-2 & Why workflow necessary & 降阶目标、审计性、不确定性和高通量处理使路线合理。 \\
Move 4-1--5 & Aim / Methods / Scope / Contributions / Structure & 本论文具体方案、边界、结果和章节组织。 \\
\bottomrule
\end{tabular}
\end{frame}

\begin{frame}{Source Keys}
\small
本 slide deck 使用论文 bibliography 中的 citation keys 作轻量标注；正式参考文献仍以 thesis PDF 的 bibliography 为准。

\vspace{0.7em}
\begin{tabular}{L{0.36\linewidth}L{0.52\linewidth}}
\toprule
主题 & source keys \\
\midrule
航运与替代燃料 & imo2023ghgstrategy; classnk2025alternativefuels \\
pilot fuel 背景 & wingd\_xdf\_manual\_2021; wingd\_lng\_faq\_2025; classnk\_methanol\_2024; wingd\_ammonia\_faq\_2025 \\
喷雾实验与模型 & Linne2013SprayImaging; HiroyasuArai1990; NaberSiebers1996; liu2020sprayreview \\
CFD 成本 & CurtisNiemeyerSung2017 \\
光学边界定义 & ecn\_liquid\_penetration; ecn\_liquid\_penetration\_accessory \\
图像处理生态 & Schindelin2012Fiji; VanDerWalt2014ScikitImage; OpenPIV; ThielickeSonntag2021PIVlab \\
\bottomrule
\end{tabular}
\end{frame}

\end{document}
